%======================================================================
% Research Track Module for SST Canon / Research Track
% Title: Core--Torsion Impedance Matching for NLSE Inertia Closure
% Author: Omar Iskandarani / Swirl-String Theory research notes
% Status: Research-track, not hard canon
%======================================================================

% Compile-safe minimal macro prelude. These commands are intentionally guarded
% so the block can be pasted into a larger Canon file that already defines them.
\providecommand{\vswirl}{\mathbf{v}_{\!\boldsymbol{\circlearrowleft}}}
\providecommand{\rc}{r_c}
\providecommand{\rhoF}{\rho_{\!f}}
\providecommand{\rhoE}{\rho_{\!E}}
\providecommand{\rhoM}{\rho_{\!m}}
\providecommand{\rhocore}{\rho_{\mathrm{core}}}
\providecommand{\GammaSST}{\Gamma_0}
\providecommand{\ct}{c_{\mathrm{T}}}
\providecommand{\clight}{c_{\mathrm{light}}}
\providecommand{\cs}{c_s}
\providecommand{\xih}{\xi}
\providecommand{\NLS}{\mathrm{NLS}}
\providecommand{\eff}{\mathrm{eff}}
\providecommand{\core}{\mathrm{core}}
\providecommand{\torsion}{\mathrm{torsion}}
\providecommand{\add}{\mathrm{add}}
\providecommand{\diag}{\mathrm{diag}}
\providecommand{\status}[1]{\textbf{[#1]}}

\section{Research Track: Core--Torsion Impedance Matching for NLSE Inertia Closure}
\label{sec:research-core-torsion-impedance}

\subsection{Purpose and epistemic status}
\label{subsec:ctim-purpose-status}

This research track records the current state of the SST inertia-closure problem after replacing the purely incompressible Euler/Biot--Savart route by a two-layer field description:
\begin{enumerate}
    \item an internal nonlinear Schrödinger/Gross--Pitaevskii-like vortex-core sector, responsible for smooth-core regularization, circulation quantization, delay-induced mode selection, and core-scale inertial response;
    \item a transverse torsion/shear radiation sector, responsible for the causal propagation speed identified with the observed light speed.
\end{enumerate}
The central point is that these two speeds should not be conflated.  The internal core sector naturally produces a characteristic acoustic speed of order the canonical swirl speed \(\lVert\vswirl\rVert\), while the transverse torsion/shear sector carries its own wave speed \(\ct\), determined by an elastic stiffness-to-density ratio.  The Lorentz-compatible clock factor can only become a derived result if the inertial response of a localized vortex core is impedance-matched to this transverse torsion sector.

\paragraph{Status declaration.}
The material in this section is \status{RESEARCH-TRACK}.  The algebraic closure
\begin{equation}
    \frac{\Delta E}{E_0}=\frac{u^2}{c_*^2}
\end{equation}
follows immediately if the same effective mass appears in both the rest-energy and translational-excitation terms.  The unresolved theorem is whether the effective mass felt by the transverse torsion sector equals
\begin{equation}
    M_{\torsion}[K] \stackrel{?}{=} \frac{2E_0[K]}{\ct^2}.
\end{equation}
Until that identity is derived from a controlled field calculation, the Lorentz clock factor remains \status{DERIVED-CONDITIONAL}, not hard canon.

\subsection{Canonical anchors and speed hierarchy}
\label{subsec:ctim-anchors}

The retained SST constants for this track are
\begin{align}
    \rc &\approx 1.40897017\times10^{-15}\ \mathrm{m},\\
    \lVert\vswirl\rVert &\approx 1.09384563\times10^{6}\ \mathrm{m\,s^{-1}},\\
    \GammaSST &= 2\pi \rc\lVert\vswirl\rVert
    \approx 9.68361920\times10^{-9}\ \mathrm{m^2\,s^{-1}},\\
    \Omega_0 &= \frac{\lVert\vswirl\rVert}{\rc}
    \approx 7.76344066\times10^{20}\ \mathrm{s^{-1}}.
\end{align}
The numerical value of \(\Omega_0\) coincides with the electron Compton angular scale in the present parameter set.  The important hierarchy is
\begin{equation}
    \lVert\vswirl\rVert \ll \clight,
    \qquad
    \frac{\lVert\vswirl\rVert}{\clight}\approx 3.64867628\times10^{-3}.
\end{equation}
Therefore any derivation that produces \(u^2/\lVert\vswirl\rVert^2\) is not yet a derivation of the Lorentz-compatible factor \(u^2/\clight^2\).  This distinction is mandatory.

\subsection{Layer I: the NLSE/Gross--Pitaevskii core sector}
\label{subsec:ctim-nlse-core}

We model the local vortex core, at research-track level, by a defocusing nonlinear Schrödinger/Gross--Pitaevskii field
\begin{equation}
    i\hbar_*\partial_t\Psi
    =
    \left[-\frac{\hbar_*^2}{2m_*}\nabla^2
    +g\left(|\Psi|^2-n_0\right)\right]\Psi,
    \label{eq:ctim-nlse}
\end{equation}
where \(m_*\) is the NLSE phase-mass parameter, \(g>0\) is the repulsive nonlinearity, and \(n_0\) is the asymptotic condensate density in the chosen normalization.  Under the Madelung representation
\begin{equation}
    \Psi=\sqrt{n}\,e^{iS},
\end{equation}
the velocity field is
\begin{equation}
    \mathbf u_{\NLS}=\frac{\hbar_*}{m_*}\nabla S,
    \label{eq:ctim-madelung-velocity}
\end{equation}
and the circulation quantum is
\begin{equation}
    \Gamma_q
    =
    \oint \mathbf u_{\NLS}\cdot d\boldsymbol\ell
    =
    \frac{h_*}{m_*}N_\Gamma,
    \qquad N_\Gamma\in\mathbb Z.
    \label{eq:ctim-circ-quantum}
\end{equation}
The associated internal sound speed and healing length are
\begin{equation}
    \cs^2=\frac{g n_0}{m_*},
    \qquad
    \xih=\frac{\hbar_*}{\sqrt{2}\,m_*\cs}.
    \label{eq:ctim-sound-healing}
\end{equation}

\paragraph{Core-speed consequence.}
If one simultaneously imposes
\begin{equation}
    \Gamma_q=\GammaSST,
    \qquad
    \xih=\rc,
\end{equation}
then Eq.~\eqref{eq:ctim-sound-healing} gives, for the standard \(h_*/m_*\) circulation normalization,
\begin{equation}
    \cs
    =
    \frac{\GammaSST}{2\sqrt{2}\pi\rc}
    =
    \frac{\lVert\vswirl\rVert}{\sqrt{2}}.
    \label{eq:ctim-cs-vswirl}
\end{equation}
Thus the NLSE core sector naturally closes on an internal acoustic speed of order \(\lVert\vswirl\rVert\), not on \(\clight\).

\paragraph{Interpretation.}
Equation~\eqref{eq:ctim-cs-vswirl} is not a defect of the model.  It says that the core healing dynamics and the electromagnetic/radiative causal dynamics are different layers.  Forcing \(\cs=\clight\) while also retaining \(\GammaSST\) and \(\xih=\rc\) overconstrains the NLSE core and breaks one of the existing SST anchors.

\subsection{Core boost inertia and the acoustic closure}
\label{subsec:ctim-core-boost}

The NLSE sector has Galilean boost covariance.  For a localized defect or vortex object \(K\), the low-velocity energy expansion may be written as
\begin{equation}
    E_{\core}[K;\mathbf u]
    =
    E_0[K]
    +
    \frac12\mathbf u^{\mathsf T}
    \mathsf M_{\core}[K]
    \mathbf u
    +
    O(\lVert\mathbf u\rVert^4),
    \label{eq:ctim-core-boost-expansion}
\end{equation}
where \(\mathsf M_{\core}[K]\) is the internal NLSE inertial tensor of the defect.  If the orientation-averaged or symmetry-reduced defect has isotropic inertial response,
\begin{equation}
    \mathsf M_{\core}[K]\to M_{\core}[K]\mathsf I,
\end{equation}
and if the rest-swirl energy satisfies the acoustic half-energy closure
\begin{equation}
    E_0[K]
    =
    \frac12M_{\core}[K]\cs^2,
    \label{eq:ctim-acoustic-half-closure}
\end{equation}
then
\begin{equation}
    \boxed{
    \frac{\Delta E_{\core}[K;u]}{E_0[K]}
    =
    \frac{u^2}{\cs^2}
    +O(u^4).
    }
    \label{eq:ctim-core-ratio}
\end{equation}
This is the correct NLSE result.  Its status is
\begin{equation}
    \boxed{\text{NLSE acoustic inertia closure: }\status{DERIVED-CONDITIONAL}.}
\end{equation}
It is not yet the Lorentz-compatible closure unless \(\cs\) is shown to equal the transverse causal speed, which current SST anchors do not support.

\subsection{Golden-NLS smooth-core branch}
\label{subsec:ctim-golden-core}

The Golden-NLS branch supplies a useful smooth-core energy regularization.  In a vortex-ring reference model one writes
\begin{align}
    E(R)
    &=
    \frac12\rho\GammaSST^2R
    \left(\ln\frac{8R}{\rc}-A\right),\\
    U(R)
    &=
    \frac{\GammaSST}{4\pi R}
    \left(\ln\frac{8R}{\rc}-B\right),
\end{align}
with core-model constants \(A\) and \(B\).  The Golden branch imposes
\begin{equation}
    B=A-1,
    \qquad
    A=\varphi,
    \qquad
    B=\varphi^{-1},
    \label{eq:ctim-golden-branch}
\end{equation}
so that
\begin{equation}
    E_{\varphi}(R)
    =
    \frac12\rho\GammaSST^2R
    \left(\ln\frac{8R}{\rc}-\varphi\right).
    \label{eq:ctim-golden-energy}
\end{equation}
This branch is useful as a reference-model bridge for \(E_0[K]\), but it should remain \status{RESEARCH-TRACK}.  It is not by itself a derivation of the transverse Lorentz speed.  Its role is to regularize the core energy and provide a controlled candidate functional for numerical comparison.

\subsection{Layer II: transverse torsion/shear causal sector}
\label{subsec:ctim-torsion-layer}

The radiative/causal layer is modeled by a transverse director or torsion/shear displacement field \(\mathbf A\) with
\begin{equation}
    \nabla\cdot\mathbf A=0.
\end{equation}
The minimal quadratic Lagrangian is
\begin{equation}
    \mathcal L_{\torsion}
    =
    \frac12\rhoF\lVert\partial_t\mathbf A\rVert^2
    -
    \frac12K\lVert\nabla\times\mathbf A\rVert^2.
    \label{eq:ctim-torsion-lagrangian}
\end{equation}
Variation gives
\begin{equation}
    \rhoF\left(\partial_t^2\mathbf A+
    \ct^2\nabla\times\nabla\times\mathbf A\right)=0,
    \qquad
    \ct^2\equiv\frac{K}{\rhoF},
    \label{eq:ctim-torsion-wave-pre}
\end{equation}
and, under \(\nabla\cdot\mathbf A=0\),
\begin{equation}
    \boxed{
    \left(\frac{1}{\ct^2}\partial_t^2-\nabla^2\right)\mathbf A=0.
    }
    \label{eq:ctim-torsion-wave}
\end{equation}
The empirical light-speed identification belongs here:
\begin{equation}
    \boxed{\ct\equiv\clight\quad\text{in the transverse radiation sector}.}
    \label{eq:ctim-c-identification}
\end{equation}
This avoids the earlier error of trying to force \(\clight\) into the core healing dynamics.

\paragraph{Torsion impedance.}
The transverse wave impedance of this sector is
\begin{equation}
    Z_{\torsion}=\rhoF\ct=\sqrt{\rhoF K}.
    \label{eq:ctim-ztorsion}
\end{equation}
The bridge problem is therefore naturally an impedance-matching problem between a localized NLSE core and the surrounding transverse torsion field.

\subsection{The Core--Torsion Impedance Matching Lemma}
\label{subsec:ctim-main-lemma}

\paragraph{Definition.}
Let \(K\) be a closed localized swirl-string configuration with internal rest-swirl energy \(E_0[K]\), internal NLSE inertial tensor \(\mathsf M_{\core}[K]\), and transverse torsion-dressed inertial tensor \(\mathsf M_{\torsion}[K]\).  The latter is defined by the quadratic low-velocity energy stored in the coupled transverse torsion field:
\begin{equation}
    \Delta E_{\torsion}[K;\mathbf u]
    =
    \frac12\mathbf u^{\mathsf T}
    \mathsf M_{\torsion}[K]
    \mathbf u
    +O(\lVert\mathbf u\rVert^4).
    \label{eq:ctim-torsion-inertia-definition}
\end{equation}

\paragraph{Research target.}
The required Lorentz-compatible closure is
\begin{equation}
    \boxed{
    \mathsf M_{\torsion}[K]
    \stackrel{?}{=}
    \frac{2E_0[K]}{\ct^2}\,\mathsf I.
    }
    \label{eq:ctim-impedance-lemma}
\end{equation}
Equivalently,
\begin{equation}
    \boxed{
    \Delta E_{\torsion}[K;u]
    \stackrel{?}{=}
    E_0[K]\frac{u^2}{\ct^2}
    +O(u^4).
    }
    \label{eq:ctim-torsion-ratio}
\end{equation}
If \(\ct=\clight\), Eq.~\eqref{eq:ctim-torsion-ratio} is the desired Lorentz-compatible energy ratio.

\paragraph{Diagnostic residual.}
Define the dimensionless closure residual
\begin{equation}
    \chi_K^{(\torsion)}
    \equiv
    \frac{\ct^2}{2E_0[K]}
    \lambda_{\iso}\!\left(\mathsf M_{\torsion}[K]\right),
    \label{eq:ctim-chi}
\end{equation}
where \(\lambda_{\iso}\) denotes the isotropic or orientation-averaged eigenvalue of the torsion-dressed inertia tensor.  The bridge closes only if
\begin{equation}
    \boxed{\chi_K^{(\torsion)}=1}
    \label{eq:ctim-chi-one}
\end{equation}
within numerical and topological tolerance.  If \(\chi_K^{(\torsion)}\neq1\) generically, then the Lorentz clock factor must remain a constitutive postulate rather than a derived theorem.

\subsection{Clock factor as a conditional consequence}
\label{subsec:ctim-clock-consequence}

If Eq.~\eqref{eq:ctim-impedance-lemma} holds, then the small-velocity energy bookkeeping becomes
\begin{equation}
    E_{\mathrm{total}}[K;u]
    =
    E_0[K]
    \left(1+\frac{u^2}{\ct^2}+O(u^4)\right)
\end{equation}
in the torsion-dressed sector.  The corresponding Lorentz-compatible clock factor may then be treated as the resummed kinetic clock relation
\begin{equation}
    \boxed{
    \frac{d\tau}{dt}
    =
    \sqrt{1-\frac{u^2}{\ct^2}}
    }
    \label{eq:ctim-clock-factor}
\end{equation}
with \(\ct=\clight\) for the transverse radiation layer.

\paragraph{Status.}
Equation~\eqref{eq:ctim-clock-factor} is not derived from the NLSE core alone.  It is \status{DERIVED-CONDITIONAL} on the impedance lemma Eq.~\eqref{eq:ctim-impedance-lemma}.  The research track therefore prevents the following incorrect inference:
\begin{equation}
    \cs=\clight.
\end{equation}
The correct statement is
\begin{equation}
    \cs\sim\lVert\vswirl\rVert
    \quad\text{for the internal core layer},
    \qquad
    \ct=\clight
    \quad\text{for the transverse torsion layer}.
\end{equation}

\subsection{Delay-induced trefoil mode selection}
\label{subsec:ctim-delay-trefoil}

The inertia-closure problem should not be separated from the topological mode-selection problem.  A closed circulation carrier has delay
\begin{equation}
    \tau_c=\frac{L}{\lVert\vswirl\rVert},
\end{equation}
and a torsion/Kelvin-like mode on the ring may be written as
\begin{equation}
    \Phi(\theta,t)=q\theta-\omega_\tau t+\phi_0.
\end{equation}
A torus-knot closure requires coprime integer windings
\begin{equation}
    p,q\in\mathbb Z,
    \qquad
    \gcd(p,q)=1.
\end{equation}
The trefoil sector is selected by
\begin{equation}
    \frac{\Omega_{\mathrm{pol}}}{\Omega_{\mathrm{tor}}}
    =
    \frac{3}{2}
    \quad\text{or}\quad
    \frac{2}{3}.
\end{equation}
This mode-locking condition constrains which circulation fields may enter the torsion bridge.  In particular, the branch label \(\sigma=\pm1\) should be interpreted as the two opposite handed circulation--torsion realizations of the same closure class, not as a literal classical spin vector.

\paragraph{Bridge implication.}
For the electron candidate \(K=T_{2,3}\), the impedance residual should be evaluated not on a generic helical packet but on the mode-locked trefoil-compatible field
\begin{equation}
    \gamma_{2{:}3}(\mathbf r,t)
    =
    \gamma(\mathbf r,t)\Big|_{q/p=3/2\ \mathrm{or}\ 2/3}.
\end{equation}
Thus the definitive numerical target is
\begin{equation}
    \chi_{T_{2,3}}^{(\torsion)}
    =
    \frac{\clight^2}{2E_0[T_{2,3}]}
    \lambda_{\iso}\!\left(\mathsf M_{\torsion}[T_{2,3};\gamma_{2{:}3}]\right)
    \stackrel{?}{=}1.
    \label{eq:ctim-trefoil-target}
\end{equation}

\subsection{Normalization issue: the \texorpdfstring{\(2\pi\)}{2pi} ambiguity}
\label{subsec:ctim-2pi}

The NLSE circulation normalization must be fixed before numerical closure tests.  Three possible mass parameters appear in the current draft ecosystem:
\begin{align}
    m_*^{(h)} &= \frac{h}{\GammaSST},\\
    m_*^{(\hbar)} &= \frac{\hbar}{\GammaSST},\\
    m_*^{(\rho V)} &= \rhocore V_{\mathrm{torus}}.
\end{align}
Only one can be the phase-mass parameter in Eq.~\eqref{eq:ctim-nlse}.  Because circulation quantization in the standard Madelung/NLSE map gives
\begin{equation}
    \Gamma_q=\frac{h}{m_*}N_\Gamma,
\end{equation}
the choice \(m_*=h/\GammaSST\) gives \(\Gamma_q=\GammaSST\) for \(N_\Gamma=1\).  If a formulation instead uses \(\hbar\), then the \(2\pi\) phase winding must be carried elsewhere.  A canon-safe implementation should declare explicitly whether the \(2\pi\) lives in the action unit, the winding number, or the definition of \(\GammaSST\).

\paragraph{Required check.}
Every numerical script in this track must report
\begin{equation}
    \mathcal R_\Gamma
    \equiv
    \frac{h/m_*}{\GammaSST}.
\end{equation}
The canonical NLSE normalization requires
\begin{equation}
    \mathcal R_\Gamma=1.
\end{equation}
If \(\mathcal R_\Gamma=2\pi\) or \(1/(2\pi)\), the calculation has not failed, but its convention must be renamed and not mixed with the canonical circulation unit.

\subsection{Minimal numerical programme}
\label{subsec:ctim-numerical-programme}

The complete research programme consists of five tests.

\paragraph{Test 1: NLSE core boost test.}
Compute a stationary vortex solution \(\Psi_K\).  Apply a small Galilean phase imprint
\begin{equation}
    \Psi_K(\mathbf r)\mapsto e^{i m_*\mathbf u\cdot\mathbf r/\hbar_*}\Psi_K(\mathbf r)
\end{equation}
and measure
\begin{equation}
    \Delta E_{\core}(u)=E_{\core}(u)-E_{\core}(0).
\end{equation}
Fit
\begin{equation}
    \Delta E_{\core}(u)=A_K^{\core}u^2+B_K^{\core}u^4+O(u^6).
\end{equation}
Then
\begin{equation}
    M_{\core}[K]=2A_K^{\core}.
\end{equation}
Compare \(M_{\core}\) to \(2E_0/\cs^2\), not to \(2E_0/\clight^2\).

\paragraph{Test 2: torsion static-dressing test.}
Treat the moving core boundary as a source for \(\mathbf A\).  Solve the transverse field equation with source
\begin{equation}
    \rhoF\left(\partial_t^2\mathbf A-\ct^2\nabla^2\mathbf A\right)
    =
    \mathbf J_K^\perp,
    \qquad
    \nabla\cdot\mathbf A=0,
\end{equation}
where \(\mathbf J_K^\perp\) is the projected current generated by the translated circulation carrier.  Compute the stored field energy
\begin{equation}
    E_{\torsion}
    =
    \frac12\rhoF\int\lVert\partial_t\mathbf A\rVert^2d^3x
    +
    \frac12K\int\lVert\nabla\times\mathbf A\rVert^2d^3x.
\end{equation}
Fit \(E_{\torsion}(u)-E_{\torsion}(0)\) to extract \(\mathsf M_{\torsion}\).

\paragraph{Test 3: impedance sweep.}
Sweep the torsion stiffness \(K\) while keeping \(\rhoF\) fixed, so that
\begin{equation}
    \ct=\sqrt{K/\rhoF}.
\end{equation}
Determine whether \(\chi_K^{(\torsion)}\) tends to unity only for \(\ct=\clight\), or whether unity requires non-canonical retuning.  A retuned-only closure is not canon.

\paragraph{Test 4: trefoil-locked versus generic-helical source.}
Repeat Test 2 for a generic helical source and for the trefoil-locked source \(\gamma_{2{:}3}\).  The bridge is topology-sensitive only if
\begin{equation}
    \chi_{T_{2,3}}^{(\torsion)}

eq
    \chi_{\mathrm{generic}}^{(\torsion)}
\end{equation}
in a stable, discretization-independent way.

\paragraph{Test 5: hierarchy consistency.}
Verify simultaneously that
\begin{equation}
    \cs\approx\frac{\lVert\vswirl\rVert}{\sqrt2},
    \qquad
    \ct=\clight,
    \qquad
    \Gamma_q=\GammaSST.
\end{equation}
If a calculation requires \(\cs=\ct\), it belongs to a different model and cannot be used as an SST canon closure.

\subsection{Falsifiers}
\label{subsec:ctim-falsifiers}

The research track fails as a derivation of Lorentz-compatible inertia if any of the following are true:
\begin{enumerate}
    \item \textbf{Non-unity torsion residual.}  The computed \(\chi_K^{(\torsion)}\) is stably different from unity for trefoil/electron-like configurations.
    \item \textbf{Non-universal topology dependence.}  The residual depends strongly on arbitrary discretization, embedding gauge, or non-topological shape choices.
    \item \textbf{Forced speed collapse.}  The calculation requires \(\cs=\ct\), contradicting the separate core and torsion speed hierarchy.
    \item \textbf{Unresolved \(2\pi\) normalization.}  The phase-mass parameter \(m_*\) fails the circulation normalization test \(h/m_* = \GammaSST\) without an explicit convention repair.
    \item \textbf{No torsion dressing.}  Integrating out the transverse torsion field contributes no quadratic inertial term tied to \(E_0/\ct^2\).
    \item \textbf{Wrong low-velocity scaling.}  The first nonzero torsion-dressed correction scales as \(u\), \(u^3\), or depends on acceleration history rather than the stored quadratic field energy.
\end{enumerate}

\subsection{Canon insertion summary}
\label{subsec:ctim-status-summary}

\begin{center}
\begin{tabular}{p{0.42\linewidth}p{0.18\linewidth}p{0.32\linewidth}}
\hline
Statement & Status & Canon interpretation \\
\hline
NLSE/GPE core sector as smooth-core reference model & \status{RESEARCH-TRACK} & Useful regularization bridge \\
\(\Gamma_q=\GammaSST,\ \xi=\rc\Rightarrow\cs=\lVert\vswirl\rVert/\sqrt2\) & \status{DERIVED-CONDITIONAL} & Internal acoustic core speed \\
\(\Delta E_{\core}/E_0=u^2/\cs^2\) & \status{DERIVED-CONDITIONAL} & Core boost closure only \\
Transverse torsion Lagrangian with \(\ct^2=K/\rhoF\) & \status{BRIDGE-LAYER} & Candidate photon/radiation sector \\
\(\ct=\clight\) & \status{CANON-CANDIDATE} & Only in transverse torsion sector \\
\(\mathsf M_{\torsion}=2E_0\ct^{-2}\mathsf I\) & \status{OPEN LEMMA} & Main research target \\
\(d\tau/dt=\sqrt{1-u^2/\clight^2}\) from this route & \status{NOT YET HARD-CANON} & Conditional on impedance lemma \\
\hline
\end{tabular}
\end{center}

\subsection{Recommended wording for the research track}
\label{subsec:ctim-recommended-wording}

The strongest defensible statement is:
\begin{quote}
SST contains two distinct continuum layers relevant to inertia closure.  The NLSE/Gross--Pitaevskii core layer supplies quantized circulation, smooth-core regularization, and an internal acoustic inertia law with speed \(\cs\sim\lVert\vswirl\rVert\).  The transverse torsion/shear layer supplies the causal radiation speed \(\ct=\clight\).  The Lorentz-compatible inertia law follows only if the localized vortex core is impedance-matched to the torsion sector such that \(\mathsf M_{\torsion}[K]=2E_0[K]\ct^{-2}\mathsf I\).  This Core--Torsion Impedance Matching Lemma is presently open and defines a falsifiable research target.
\end{quote}

\subsection{Bibliography for this research-track module}
\label{subsec:ctim-bibliography}

\begin{thebibliography}{99}

\bibitem{Helmholtz1858CTIM}
H.~von Helmholtz,
``Über Integrale der hydrodynamischen Gleichungen, welche den Wirbelbewegungen entsprechen,''
\emph{Journal für die reine und angewandte Mathematik} \textbf{55}, 25--55 (1858).

\bibitem{Kelvin1869CTIM}
W.~Thomson (Lord Kelvin),
``On Vortex Motion,''
\emph{Transactions of the Royal Society of Edinburgh} \textbf{25}, 217--260 (1869).

\bibitem{Madelung1927CTIM}
E.~Madelung,
``Quantentheorie in hydrodynamischer Form,''
\emph{Zeitschrift für Physik} \textbf{40}, 322--326 (1927).
DOI: 10.1007/BF01400372.

\bibitem{Pitaevskii1961CTIM}
L.~P. Pitaevskii,
``Vortex Lines in an Imperfect Bose Gas,''
\emph{Soviet Physics JETP} \textbf{13}, 451--454 (1961).

\bibitem{Gross1961CTIM}
E.~P. Gross,
``Structure of a Quantized Vortex in Boson Systems,''
\emph{Il Nuovo Cimento} \textbf{20}, 454--477 (1961).
DOI: 10.1007/BF02731494.

\bibitem{PitaevskiiStringari2003CTIM}
L.~Pitaevskii and S.~Stringari,
\emph{Bose-Einstein Condensation},
Oxford University Press, 2003.

\bibitem{Moffatt1969CTIM}
H.~K. Moffatt,
``The degree of knottedness of tangled vortex lines,''
\emph{Journal of Fluid Mechanics} \textbf{35}, 117--129 (1969).
DOI: 10.1017/S0022112069000991.

\bibitem{Batchelor1967CTIM}
G.~K. Batchelor,
\emph{An Introduction to Fluid Dynamics},
Cambridge University Press, 1967.

\bibitem{Saffman1992CTIM}
P.~G. Saffman,
\emph{Vortex Dynamics},
Cambridge University Press, 1992.

\bibitem{BerryDennis2001CTIM}
M.~V. Berry and M.~R. Dennis,
``Knotted and linked phase singularities in monochromatic waves,''
\emph{Proceedings of the Royal Society A} \textbf{457}, 2251--2263 (2001).
DOI: 10.1098/rspa.2001.0826.

\bibitem{Kedia2013CTIM}
H.~Kedia, I.~Bialynicki-Birula, D.~Peralta-Salas, and W.~T.~M. Irvine,
``Tying knots in light fields,''
\emph{Physical Review Letters} \textbf{111}, 150404 (2013).
DOI: 10.1103/PhysRevLett.111.150404.

\bibitem{BarceloLiberatiVisser2011CTIM}
C.~Barceló, S.~Liberati, and M.~Visser,
``Analogue Gravity,''
\emph{Living Reviews in Relativity} \textbf{14}, 3 (2011).
DOI: 10.12942/lrr-2011-3.

\bibitem{Volovik2003CTIM}
G.~E. Volovik,
\emph{The Universe in a Helium Droplet},
Oxford University Press, 2003.

\end{thebibliography}
